% Бүлэг 4
\chapter{Дүгнэлт}
\label{Chapter4}

\section{Дүгнэлт}
\qquad
Энэхүү бие даалтын ажлын хүрээнд \textbf{Эмийн бүртгэлийн систем}-ийг боловсруулан нэвтрүүлэх боломжийг судалж, эмийн мэдээллийг цахим хэлбэрээр бүртгэх, удирдах, хянах автоматжуулсан шийдлийг боловсруулах зорилгыг тавьсан. Судалгааны үр дүнд тус систем нь эмийн сангийн үйл ажиллагааг хялбарчлах, мэдээллийн алдааг бууруулах, ажлын бүтээмжийг нэмэгдүүлэхэд чухал ач холбогдолтой болох нь тогтоогдлоо.

Эмийн бүртгэлийн системийг нэвтрүүлснээр дараах үр өгөөжийг бий болгоно:
\begin{enumerate}
\item \textbf{Ажлын үр ашиг нэмэгдэх} – Эмийн бүртгэл, захиалга, нийлүүлэлт, үлдэгдлийн мэдээллийг автоматжуулснаар гараар бүртгэл хийх хэрэгцээ багасч, хүний алдааг бууруулна.
\item \textbf{Нөөцийн хяналт сайжрах} – Эмийн үлдэгдэл, хугацаа дуусах дөхсөн эмийг бодит цаг хугацаанд хянах боломж бүрдэж, хомсдол болон хэт их нөөц үүсэхээс сэргийлнэ.
\item \textbf{Цаг хугацаа хэмнэх} – Захиалга, борлуулалт, тайлан гаргалтыг автоматжуулснаар ажилтнуудын ачаалал буурч, үйл ажиллагааны хурд нэмэгдэнэ.
\item \textbf{Мэдээллийн ил тод байдал} – Бүх мэдээлэл төвлөрсөн сангаас харагдах боломжтой болж, удирдлагын шийдвэр гаргалт илүү шуурхай, оновчтой болно.
\item \textbf{Үйлчлүүлэгчийн сэтгэл ханамж нэмэгдэх} – Эмийн мэдээлэл шуурхай, үнэн зөв хүргэгдэх, захиалга хурдан боловсруулагдах зэргээр үйлчилгээний чанар сайжирна.
\end{enumerate}

Эмийн бүртгэлийн систем нь орчин үеийн технологийн шаардлагад нийцсэн, эмийн сангийн үйл ажиллагааг оновчтой зохион байгуулах чухал хэрэгсэл юм. Цаашид уг системийг илүү өргөн хүрээнд хөгжүүлж, онлайн захиалга, санхүүгийн модуль, хиймэл оюун ухаанд суурилсан нөөцийн таамаглалыг нэмснээр эмийн сангийн салбарт илүү үр ашигтайгаар ашиглах бүрэн боломжтой гэж үзэж байна.

\newpage

\section{Ашигласан материалын жагсаалт}

\begin{center}
Ном, судалгааны материал
\end{center}

\begin{center}
Вэб сайтууд
\end{center}
\begin{enumerate}
\item https://www.legalinfo.mn/law/details/12924?lawid=12924
\item https://www.who.int
\item https://www.google.mn/
\item https://www.tutorialspoint.com/
\item https://stackoverflow.com/
\end{enumerate}
